\section{Discussion}\label{sec:discussion}

\textbf{What the NO-GO means.} The decision rule asks for evidence that a
choice per document beats the best single setup, and neither held-out set
gives it. This is not proof that a better chooser can never help: the best
single setup is selected on the same questions it is scored on, and the
oracle chooses from only half of a document's few questions. The result is
better read as follows: in this design, many setups are statistically tied,
the differences between them are small, and selecting on a handful of
questions per document adds more noise than it removes. The larger choice
sets point the same way: with more setups to choose from, the oracle does
worse, not better. In practice, this supports a simpler default: use one
small-chunk setup and spend tuning effort on its chunk size rather than on a
per-document selector.

\textbf{The measure decides the answer.} RQ2 and RQ3 together show that the
best chunk size is not a property of the chunker alone. Under a fixed number
of chunks, a larger chunk shows the reader more text, so Recall@$k$ rewards
size itself; under a fixed token budget that advantage disappears and small
chunks tend to win. Comparisons of chunk sizes that fix the number of
retrieved chunks therefore mix the quality of the chunks with the amount of
text read. We suggest that such comparisons report the token budget, or both
measures. Our size effect also weakens as the budget grows and is not
established at 2,000 tokens on held-out set 2; a larger budget leaves room
for larger chunks, but we did not test this explanation.

\textbf{Defaults and tuning.} Tuning paid off in one place: LangChain's
default chunk size was far from the best size for this budget. Where our
starting size was already close, tuning brought no detectable gain, and
parser settings changed recall very little. For scanned PDFs, the parser
choice that mattered was whether text was extracted at all. A practical check
is therefore to compare a library's default chunk size with the token budget
the system will use, and to detect scanned pages before parsing.

\textbf{Answers and retrieval.} Answer quality agrees with retrieval on the
large-chunk penalty, but the evidence is weak: it rests on few setups, one
answering model and a measure chosen after the answers existed. Whether
retrieval gains of the size we measure change answers in practice needs a
larger, pre-registered answer study.

\textbf{Future work.} A chooser could still help with more questions per
document, other selection signals, other languages or more varied parsers. A
fair test of such a chooser should keep the conditions used here: a fixed
token budget, sealed held-out documents, and an oracle that cannot profit from
selection luck.
